\section{Zadanie 3}
\autor{Dominik Korsa}

Udowodnij, że dla dowolnego ciągu \(A_0, A_1\), ... zachodzi:

\[
	P(\bigcup_{n = 0}^{\infty} A_n) \leq \sum_{n = 0}^{\infty} P(A_n)
\]

\subsection*{Rozwiązanie}
Definiujemy nowy ciąg \textbf{rozłącznych zdarzeń}:
\(\displaystyle
B_n = A_n \setminus \bigcup_{i = 0}^{n-1} A_i
\).

Można pokazać, że:
\(\displaystyle
\bigcup_{n \in \mathbb{N}} B_n = \bigcup_{n \in \mathbb{N}} A_n
\).

\begin{align*}
	P\left(\bigcup_{n \in \mathbb{N}} A_n\right)
	 & = P\left(\bigcup_{n \in \mathbb{N}} B_n\right) \\
	 & = \sum_{n \in \mathbb{N}} P(B_n)               \\
	 & \leq \sum_{n \in \mathbb{N}} P(A_n)
\end{align*}
Pierwsze przejście mamy z definicji miary probabilistycznej, dzięki temu że
zdarzenia w rodzinie \(\{B_n\}\) są parami rozłączne.

Nierówność mamy, bo dla każdego \(n \in \mathbb{N}\) mamy \(B_n \subseteq A_n\), więc \(P(B_n) \leq P(A_n)\).

Otrzymana nierówność kończy dowód.
